\chapter{Estado de la Cuestión}
\label{cap:estadoDeLaCuestion}

Durante décadas, la evolución de la computación se apoyó en aumentar la frecuencia de los
procesadores y en integrar más transistores en el mismo chip. Esa tendencia permitió que muchas
aplicaciones mejorasen sin cambiar demasiado la forma de programarlas. Sin embargo, en los últimos
años se ha hecho evidente que esa vía no escala igual de bien: el consumo, la disipación térmica y
la dificultad de explotar más paralelismo dentro de una CPU generalista han empujado a buscar
arquitecturas más especializadas \citep{10.1145/3282307}.

Este cambio se nota especialmente en aplicaciones de inteligencia artificial. Una red neuronal no
suele estar limitada por una única operación compleja, sino por miles o millones de operaciones
simples repetidas: multiplicaciones, acumulaciones, accesos a memoria, activaciones y movimientos
de datos. En una CPU estas operaciones se ejecutan siguiendo una arquitectura muy flexible, pero no
siempre eficiente. En cambio, una FPGA permite construir una ruta de datos a medida para una tarea
concreta, explotando paralelismo espacial y reutilizando recursos de forma controlada
\citep{electronics10091025}.

Por este motivo, muchos sistemas empotrados actuales se plantean como sistemas heterogéneos: una
parte del chip se dedica al control y a la ejecución de software convencional, mientras que otra
parte implementa aceleradores específicos. En este trabajo esa división aparece de forma natural:
el procesador ARM prepara los datos, controla las transferencias y muestra los resultados, mientras
que la lógica programable ejecuta la inferencia del modelo.

\section{Arquitectura Zynq-7000 y Comunicaciones AXI}
La familia Zynq-7000 combina en un mismo circuito integrado un sistema de procesamiento
(\textit{Processing System}, PS) y lógica programable (\textit{Programmable Logic}, PL). En el caso
de la Zybo utilizada en este trabajo, el dispositivo es un Zynq Z-7010. La parte PS incluye un ARM
Cortex-A9 de doble núcleo, memoria DDR y periféricos habituales de un sistema empotrado. La parte
PL equivale a una FPGA de la familia 7-series, donde se pueden integrar bloques hardware propios.

Esta arquitectura es especialmente interesante porque evita tener una FPGA aislada. El procesador y
la lógica programable comparten el mismo sistema y se comunican mediante buses AMBA AXI
\citep{amdAxiDma}. AXI4-Lite
se utiliza normalmente para acceder a registros de control, por ejemplo para arrancar un IP o leer
su estado. AXI4-Stream, en cambio, está pensado para transportar datos como un flujo continuo, sin
direcciones por cada palabra. Por último, los puertos de alto rendimiento del Zynq permiten que un
DMA lea y escriba directamente en memoria DDR, descargando al procesador de copiar los datos uno a
uno.

Esta separación entre control y datos es clave en aceleración. El ARM no debe enviar cada pixel de
forma manual al acelerador, porque entonces el coste de comunicación consumiría buena parte del
beneficio. La solución habitual es preparar buffers en DDR, limpiar o invalidar la caché cuando sea
necesario, y dejar que el AXI DMA mueva bloques completos entre memoria y FPGA.

La limitación principal de la Zybo no está en el concepto de arquitectura, sino en el tamaño del
dispositivo. Según la hoja de datos de la familia Zynq-7000, el Zynq Z-7010 dispone de 17.600 LUT,
35.200 flip-flops, 60 bloques BRAM de 36 Kb y 80 bloques DSP \citep{zynq7000Datasheet}. Estas
cifras son suficientes para aprender y prototipar, pero obligan a cuidar mucho el tamaño de una red
convolucional.

\section{Síntesis de Alto Nivel (HLS)}
Tradicionalmente, una FPGA se programaba describiendo el hardware con VHDL o Verilog. Esto da mucho
control, pero también obliga a pensar en registros, ciclos de reloj, máquinas de estados y rutas de
datos desde el principio. La Síntesis de Alto Nivel, o HLS, intenta elevar ese nivel de abstracción:
el diseñador describe el algoritmo en \texttt{C} o \texttt{C++}, y la herramienta genera
una implementación RTL que
después puede integrarse como IP en Vivado.

En el caso de AMD/Xilinx, Vitis HLS sintetiza funciones \texttt{C}/\texttt{C++} a RTL
para implementarlas en la
lógica programable y permite exportarlas como IP de Vivado \citep{amdVitisHLS}. Esto no significa
que el código C se ejecute como software dentro de la FPGA. La herramienta analiza bucles, arrays,
dependencias de datos y tipos, y a partir de ahí construye hardware: memorias, registros,
multiplexores, operadores aritméticos y lógica de control.

Las directivas HLS son una parte importante del flujo. Con \texttt{PIPELINE} se puede intentar que
un bucle acepte nuevos datos cada pocos ciclos, aunque las operaciones anteriores sigan avanzando
por la ruta de datos. Con \texttt{UNROLL} se pueden duplicar operaciones para ejecutar varias
iteraciones en paralelo. Con directivas de interfaz se decide si una función se conectará por
AXI4-Lite, AXI4-Stream o memoria. Estas decisiones no son detalles menores: afectan directamente a
la latencia, al consumo de LUT, al uso de DSP, a la BRAM y a la posibilidad de que Vivado coloque
el diseño en la FPGA.

Por eso HLS no elimina el diseño hardware, sino que cambia dónde se toman muchas decisiones. En un
proyecto pequeño puede acelerar mucho el desarrollo; en un proyecto ajustado en recursos, como una
CNN en una Zynq Z-7010, sigue siendo imprescindible leer los informes de síntesis, entender los
cuellos de botella y ajustar el código o las directivas.

\section{hls4ml y generación de modelos para FPGA}
hls4ml \citep{fastml_hls4ml} es una herramienta de codiseño hardware-software orientada a traducir
modelos de aprendizaje
automático a implementaciones para FPGA o ASIC. En lugar de escribir a mano cada capa de una red
neuronal en \texttt{C++} para HLS, hls4ml parte de modelos entrenados en herramientas
habituales de machine
learning y genera un proyecto HLS con la topología, los pesos y las funciones necesarias para la
inferencia \citep{fahim2021hls4mlopensourcecodesignworkflow}.

Su flujo encaja muy bien con este proyecto. Primero se entrena un modelo en TensorFlow/Keras y se
guarda en formato \texttt{.h5}. Después hls4ml interpreta esa red, crea una representación interna
del grafo y genera código \texttt{C++} sintetizable. La documentación de hls4ml describe
esta separación
entre frontends, que leen modelos de frameworks como Keras u ONNX, y backends, que generan código
para herramientas como Vivado HLS, Vitis HLS o Intel HLS \citep{hls4mlDocs}.

La aportación de hls4ml no es solo convertir capas. También permite configurar la precisión de los
datos, el factor de reutilización de operadores, la estrategia de síntesis y el tipo de entrada y
salida del modelo. En una FPGA pequeña, estas opciones son decisivas. Reducir el ancho de palabra
puede ahorrar LUT y BRAM, pero también puede degradar la precisión. Aumentar el factor de
reutilización reduce área, pero aumenta latencia. Elegir una arquitectura de flujo puede ahorrar
memoria en algunas capas, pero introduce FIFO y lógica adicional que también cuentan para Vivado.

En este trabajo se ha usado hls4ml precisamente como puente entre el entrenamiento del modelo y la
generación del IP. La red no se ha tratado como una caja negra: ha sido necesario ajustar canales,
precisiones, buffers, interfaz AXI-Stream y configuración de recursos hasta conseguir que el diseño
tuviera un tamaño apto para la Zybo.

\section{IA en FPGA con recursos limitados}
El despliegue de modelos de visión artificial en sistemas empotrados está condicionado por dos
restricciones: cálculo y memoria. En una GPU o en una FPGA grande se puede ejecutar una red
convolucional con muchas capas y canales, pero en una Zynq Z-7010 cada capa adicional se nota en el
uso de LUT, BRAM y lógica de control.

Por eso, cuando se trabaja con placas pequeñas, se suelen aplicar modelos compactos y técnicas de
cuantización. La cuantización reduce la precisión de pesos y activaciones para que la inferencia
pueda ejecutarse con enteros o punto fijo en lugar de coma flotante. Esta idea está ampliamente
extendida en inferencia embebida: trabajos como el de Jacob et al. muestran que la aritmética
entera puede reducir memoria y mejorar la eficiencia manteniendo una precisión razonable si se
diseña correctamente el flujo de entrenamiento e inferencia \citep{Jacob_2018_CVPR}.

En este proyecto se ha seguido esa misma filosofía, aunque adaptada al flujo HLS. El modelo se
entrena en coma flotante, pero la inferencia final se ejecuta con tipos de punto fijo
\citep{amdPrecisionFixedPoint}. Esta
decisión permite que el diseño sea viable en recursos, pero introduce una consecuencia importante:
el comportamiento numérico de la FPGA no tiene por qué coincidir exactamente con el modelo original
de Keras. Los redondeos, truncados y anchos de acumulación forman parte del diseño.

CIFAR-10 añade además una dificultad razonable para una placa de este tamaño. El dataset contiene
60.000 imágenes RGB de 32x32 repartidas en 10 clases, con 50.000 imágenes de entrenamiento y
10.000 de test \citep{cifar10Dataset}. Frente a MNIST, que trabaja con dígitos en escala de grises,
CIFAR-10 exige extraer características de color, textura y forma en imágenes muy pequeñas. Esto lo
convierte en una prueba más interesante, pero también más exigente para una CNN reducida.

\section{Herramientas utilizadas en el flujo}
El proyecto se apoya en varias herramientas que cubren etapas distintas del desarrollo:

\begin{itemize}
  \item \textbf{TensorFlow/Keras:} entrenamiento del modelo y exportación a formato \texttt{.h5}.
  \item \textbf{hls4ml:} conversión del modelo entrenado a código \texttt{C++}
    sintetizable y generación del
    proyecto HLS.
  \item \textbf{Vitis HLS:} síntesis del código \texttt{C++} a RTL y empaquetado del
    acelerador como IP.
  \item \textbf{Vivado:} integración del IP con el Zynq Processing System, AXI DMA e interconexiones
    AXI, además de síntesis, implementación y generación del bitstream.
  \item \textbf{Vitis:} desarrollo de las aplicaciones bare-metal que se ejecutan sobre el ARM,
    inicializan los drivers, preparan buffers y lanzan la inferencia.
\end{itemize}

Esta cadena de herramientas permite recorrer todo el camino desde el entrenamiento hasta la placa,
pero también introduce un reto práctico: cada etapa genera artefactos intermedios y puede conservar
resultados antiguos. Por eso la reproducibilidad y la organización de carpetas son parte importante
del trabajo, no solo una cuestión estética.

\section{Panorama industrial y proyectos de referencia}
La transición hacia la computación heterogénea no es solo una respuesta académica al estancamiento
del rendimiento de las CPU. También se ha convertido en una estrategia industrial para reducir
latencia y consumo en cargas muy concretas.

En conjunto, el estado de la cuestión muestra que el proyecto se sitúa en una línea muy actual: la
necesidad de mover inferencia de IA hacia hardware especializado y cercano al dato. La Zybo no puede
competir con aceleradores modernos, pero sí permite estudiar a pequeña escala los mismos problemas:
movimiento de datos, cuantización, generación de IPs, integración AXI y equilibrio entre precisión,
latencia y recursos.

\subsection{Microsoft Project Catapult}
Uno de los referentes más conocidos en el uso de FPGA a gran escala es Project Catapult de
Microsoft\footnote{\url{https://www.microsoft.com/en-us/research/project/project-catapult/}}.
Este proyecto surgió para acelerar tareas de búsqueda y
clasificación en centros de datos \citep{putnam2014a}.
La idea principal era añadir FPGA a los servidores para obtener aceleradores reconfigurables
capaces de adaptarse a distintos algoritmos sin sustituir toda la infraestructura.

\subsection{NVDLA}
NVDLA\footnote{\url{https://nvdla.org/}}, anunciado por NVIDIA, representa otro enfoque:
una arquitectura abierta y configurable para
inferencia de aprendizaje profundo \citep{nvdlaDocumentation}.
Aunque está más cerca de un acelerador dedicado que de una FPGA
general, es relevante porque refleja la misma tendencia: construir hardware específico para redes
neuronales, con especial atención a tipos de datos reducidos y eficiencia.

\subsection{BNN-PYNQ-ZYBO}
A una escala más cercana a este trabajo, el repositorio
\textit{BNN-PYNQ-ZYBO}\footnote{\url{https://github.com/fcaspe/BNN-PYNQ-ZYBO}} adapta el proyecto
BNN-PYNQ \citep{finn} para la Zybo-Z7-10. No es exactamente la Zybo Legacy utilizada aquí, pero
sí trabaja con
una placa basada en el mismo tamaño de dispositivo, el Zynq Z-7010. El propio repositorio indica que
incluye topologías pensadas específicamente para la Zybo-Z7-10, con redes de pesos y
activaciones de 1 bit, scripts de síntesis para esa placa y soporte para conjuntos como CIFAR-10 y
MNIST.

Este tipo de proyecto es útil como referencia porque confirma una idea que también aparece durante
este TFG: en un Zynq Z-7010, la clave no es intentar llevar una red grande sin cambios, sino diseñar
una red muy compacta, cuantizada y ajustada al hardware disponible. La diferencia es que este
trabajo no parte de una red binaria ya preparada, sino que recorre el flujo completo desde
TensorFlow/Keras hasta hls4ml, Vitis HLS, Vivado y la aplicación bare-metal.
